\documentclass{exam-zh}
\usepackage{siunitx}
\usepackage[most]{tcolorbox}

\examsetup{
  page/size=a4paper,
  paren/show-paren=true,
  paren/show-answer=true,
  fillin/show-answer=false,
  solution/show-solution=show-stay,
  solution/label-indentation=false
}

\everymath{\displaystyle}

\title{{{TITLE}}}
% \subject{{{SUBJECT}}}

\begin{document}

\information{
  姓名\underline{\hspace{6em}},
  学号\underline{\hspace{6em}},
  班级\underline{\hspace{6em}}
}

% \secret

\maketitle

% \begin{notice}
%   \item 答卷前，考生务必将自己的姓名、考生号、考场号和座位号填写答题卡上，用2B铅笔将试卷类型（B）填涂在答题卡相应位置上，将条形码横贴在答题卡右上角“条形码粘贴处”。
%   \item 作答选择题时，选出每小题答案后，用2B铅笔在答题卡上对应题目选项的答案信息点涂黑；如需改动，用橡皮擦干净后，再选涂其他答案。答案不能答在试卷上。
%     写在试卷、草稿纸和答题卡上的非答题区域均无效。
%   \item 非选择题必须用黑色字迹的钢笔或签字笔作答，答案必须写在答题卡各题目指定区域内相应位置上；如需改动，先划掉原来的答案，然后再写上新答案；不准使用铅笔和涂改液。不按以上要求作答无效。
%   \item 考生必须保持答题卡的整洁。考试结束后，将试卷和答题卡一并交回。
% \end{notice}

{{QUESTION_SECTIONS}}

\end{document}
